% answers_short.tex —— 2009 年·春季上海卷 · 简短答案速查
% 由 tools/make_short_answers.py 自 answers.tex 的 \gksolhead/\gkanswer 提取，
% 并与 2009/2009_数学真题_春季上海_答案要点.md 逐题比对后生成，勿手改（改源 MD 后重跑）。

\begin{gknotice}
\textbf{使用说明}：本节只给最终答案便于快速核对，完整推导见解析版。答案为编写者独立推导，\textbf{非官方标准答案}。
\end{gknotice}

\gkpart{填空题}
\begin{itemize}[leftmargin=2.6em]
  \item[\textbf{1.}] $(1,+\infty)$
  \item[\textbf{2.}] $-2i$
  \item[\textbf{3.}] $4\pi$
  \item[\textbf{4.}] $(1,2)$
  \item[\textbf{5.}] $x=-\dfrac14$
  \item[\textbf{6.}] $120^\circ$
  \item[\textbf{7.}] $y=x-5$
  \item[\textbf{8.}] $\sqrt6$
  \item[\textbf{9.}] $0$
  \item[\textbf{10.}] $\dfrac{5}{26^6}$
  \item[\textbf{11.}] 第二次为 $\{1/4,3/4\}$；第 $n$ 次为 $\{(2k-1)/2^n:1\le k\le2^{n-1}\}$
\end{itemize}

\gkpart{选择题}
\begin{center}\begin{tabular}{ccccc}
\toprule
题号 & 12 & 13 & 14 & 15 \\
\midrule
答案 & B & C & A & C \\
\bottomrule
\end{tabular}\end{center}

\gkpart{解答题}
\begin{itemize}[leftmargin=2.6em]
  \item[\textbf{16.}] U（现存题面条件不足以唯一确定体积）
  \item[\textbf{17.}] (1) $a_n=3^{-(n-1)}$；(2) $k_{\max}=\dfrac23$
  \item[\textbf{18.}] 约 $187$ 百公里，即 $18{,}700$ km
  \item[\textbf{19.}] (1) $A_1,A_2,A_3$ 不共线；(2) 所有 $A_n$ 均满足 $y_n^2=2x_n$；(3) $a=2d$
  \item[\textbf{20.}] (1) $f_1,f_3$ 均递增；(2) 题设恒等式成立；(3) 奇数 $n$ 的最大值为 $0$、最小值为 $-1$；偶数 $n$ 的最大值为 $1$、最小值为 $2^{1-n/2}$
\end{itemize}
